
\documentclass{article}
\usepackage{colortbl}
\usepackage{makecell}
\usepackage{multirow}
\usepackage{supertabular}

\begin{document}

\newcounter{utterance}

\centering \large Interaction Transcript for game `cladder', experiment `full\_v1.5\_default', episode 9295 with reflexion{-}Qwen3.5{-}9B.
\vspace{24pt}

{ \footnotesize  \setcounter{utterance}{1}
\setlength{\tabcolsep}{0pt}
\begin{supertabular}{c@{$\;$}|p{.15\linewidth}@{}p{.15\linewidth}p{.15\linewidth}p{.15\linewidth}p{.15\linewidth}p{.15\linewidth}}
    \# & \multicolumn{2}{c}{Player} && \multicolumn{2}{c}{Game Master} \\
    \hline

    \theutterance \stepcounter{utterance}  
    & & & \multicolumn{4}{p{0.6\linewidth}}{
        \cellcolor[rgb]{0.9,0.9,0.9}{
            \makecell[{{p{\linewidth}}}]{
                \texttt{\tiny{[P1$\langle$GM]}}
                \texttt{You have attempted this same task before and did not succeed. These are the plans you wrote after those attempts.} \\
\texttt{Trial \#1: New plan: I failed because I stopped my reasoning mid{-}calculation after determining the numerator for the conditional probability, never completing the division to find \$P(\textbackslash{}text\{glimx\} \textbackslash{}mid \textbackslash{}text\{zuph\})\$ or comparing it to the marginal \$P(\textbackslash{}text\{glimx\})\$. Specifically, I calculated \$P(\textbackslash{}text\{glimx\}) = 0.31\$ but did not finish calculating \$P(\textbackslash{}text\{glimx\} \textbackslash{}mid \textbackslash{}text\{zuph\}) = \textbackslash{}frac\{0.30\}\{0.96\} \textbackslash{}approx 0.3125\$. Since \$0.3125 > 0.31\$, the chance is actually larger, not smaller. Next time, I will explicitly calculate the full value of the conditional probability, compare it numerically to the marginal probability, and only then conclude with "Yes" or "No" based on that comparison.} \\
\\ 
\texttt{Trial \#2: New plan: I failed because I incorrectly calculated the marginal probability \$P(\textbackslash{}text\{glimx\})\$ as \$0.31\$ instead of \$0.33\$ (summing \$0.30\$ and \$0.03\$). This arithmetic error led me to conclude that \$0.3125 < 0.31\$, whereas the correct comparison is \$0.3125 < 0.33\$, which actually supports the "Yes" answer. Next time, I will explicitly verify the sum of disjoint probabilities before comparing them to the conditional probability, ensuring \$P(\textbackslash{}text\{glimx\}) = 0.30 + 0.03 = 0.33\$. Since \$0.3125 < 0.33\$, the correct conclusion is "Yes".} \\
\\ 
\texttt{Trial \#3: New plan: I failed because I incorrectly calculated the marginal probability \$P(\textbackslash{}text\{glimx\})\$ as \$0.31\$ instead of the correct sum \$0.33\$ (\$0.30 + 0.03\$). This arithmetic error led me to believe the conditional probability (\$0.3125\$) was larger than the marginal, resulting in a "No" answer. Next time, I will explicitly verify the sum of disjoint probabilities first: \$P(\textbackslash{}text\{glimx\}) = P(\textbackslash{}text\{zuph\} \textbackslash{}cap \textbackslash{}text\{glimx\}) + P(\textbackslash{}neg\textbackslash{}text\{zuph\} \textbackslash{}cap \textbackslash{}text\{glimx\}) = 0.30 + 0.03 = 0.33\$. Then, comparing the correct values (\$0.3125 < 0.33\$), I will conclude "Yes" because the chance of glimx is indeed smaller when observing zuph.} \\
\\ 
\\ 
\texttt{You are an expert in causal inference. The following question is not a typical commonsense query, but rather a meticulously designed question created by a professor specializing in causal inference, intended to assess the students' mastery of the course content.} \\
\texttt{Imagine a self{-}contained, hypothetical world with only the following conditions, and without any unmentioned factors or causal relationships: Zory has a direct effect on zuph and glimx. Zuph has a direct effect on jyka. Jyka has a direct effect on glimx. Zory is unobserved. The overall probability of zuph is 96\%. The probability of not zuph and glimx is 3\%. The probability of zuph and glimx is 30\%. Is the chance of glimx smaller when observing zuph?} \\
\texttt{Start your answer with 'Yes' or 'No', followed by additional reasoning or evidence to support your explanation.} \\
            }
        }
    }
     \\ \\

    \theutterance \stepcounter{utterance}  
    & \multicolumn{4}{p{0.6\linewidth}}{
        \cellcolor[rgb]{0.561,0.737,0.561}{
            \makecell[{{p{\linewidth}}}]{
                \texttt{\tiny{[P1$\rangle$GM]}}
                \texttt{Yes} \\
\\ 
\texttt{**Reasoning:**} \\
\\ 
\texttt{To determine if the chance of `glimx` is smaller when observing `zuph`, we must compare the conditional probability \$P(\textbackslash{}text\{glimx\} \textbackslash{}mid \textbackslash{}text\{zuph\})\$ with the marginal probability \$P(\textbackslash{}text\{glimx\})\$.} \\
\\ 
\texttt{1.  **Calculate the Marginal Probability \$P(\textbackslash{}text\{glimx\})\$:**} \\
\texttt{    The problem provides the probabilities for the two mutually exclusive states of `glimx`:} \\
\texttt{    *   \$P(\textbackslash{}text\{zuph\} \textbackslash{}cap \textbackslash{}text\{glimx\}) = 0.30\$} \\
\texttt{    *   \$P(\textbackslash{}neg\textbackslash{}text\{zuph\} \textbackslash{}cap \textbackslash{}text\{glimx\}) = 0.03\$ (given as "probability of not zuph and glimx")} \\
\\ 
\texttt{    Summing these gives the total probability of `glimx`:} \\
\texttt{    \$\$P(\textbackslash{}text\{glimx\}) = 0.30 + 0.03 = 0.33\$\$} \\
\\ 
\texttt{2.  **Calculate the Conditional Probability \$P(\textbackslash{}text\{glimx\} \textbackslash{}mid \textbackslash{}text\{zuph\})\$:**} \\
\texttt{    The definition of conditional probability is \$P(A \textbackslash{}mid B) = \textbackslash{}frac\{P(A \textbackslash{}cap B)\}\{P(B)\}\$.} \\
\texttt{    *   The numerator is the joint probability \$P(\textbackslash{}text\{zuph\} \textbackslash{}cap \textbackslash{}text\{glimx\}) = 0.30\$.} \\
\texttt{    *   The denominator is the marginal probability of the condition, \$P(\textbackslash{}text\{zuph\})\$. The problem states the overall probability of `zuph` is \$96\textbackslash{}\%\$, or \$0.96\$.} \\
\\ 
\texttt{    \$\$P(\textbackslash{}text\{glimx\} \textbackslash{}mid \textbackslash{}text\{zuph\}) = \textbackslash{}frac\{0.30\}\{0.96\} = 0.3125\$\$} \\
\\ 
\texttt{3.  **Compare the Probabilities:**} \\
\texttt{    *   \$P(\textbackslash{}text\{glimx\} \textbackslash{}mid \textbackslash{}text\{zuph\}) = 0.3125\$} \\
\texttt{    *   \$P(\textbackslash{}text\{glimx\}) = 0.33\$} \\
\\ 
\texttt{    Since \$0.3125 < 0.33\$, the probability of `glimx` occurring is indeed lower when `zuph` is observed compared to the overall average.} \\
\\ 
\texttt{Therefore, the answer is **Yes**.} \\
            }
        }
    }
    & & \\ \\

    \theutterance \stepcounter{utterance}  
    & & & \multicolumn{2}{p{0.3\linewidth}}{
        \cellcolor[rgb]{0.9,0.9,0.9}{
            \makecell[{{p{\linewidth}}}]{
                \texttt{\tiny{[GM$|$GM]}}
                \texttt{Yes} \\
            }
        }
    }
    & & \\ \\

    \theutterance \stepcounter{utterance}  
    & & & \multicolumn{2}{p{0.3\linewidth}}{
        \cellcolor[rgb]{0.9,0.9,0.9}{
            \makecell[{{p{\linewidth}}}]{
                \texttt{\tiny{[GM$|$GM]}}
                \texttt{yes} \\
            }
        }
    }
    & & \\ \\

    \theutterance \stepcounter{utterance}  
    & & & \multicolumn{2}{p{0.3\linewidth}}{
        \cellcolor[rgb]{0.9,0.9,0.9}{
            \makecell[{{p{\linewidth}}}]{
                \texttt{\tiny{[GM$|$GM]}}
                \texttt{game\_result = WIN} \\
            }
        }
    }
    & & \\ \\

\end{supertabular}
}

\end{document}
